%=============================================================
\chapter{Implementação serial}
\label{cap:serial}
%=============================================================

A implementação serial do modelo, escrita na linguagem C, constitui a referência deste trabalho, uma vez que é contra os seus resultados que a implementação paralela será validada. Neste capítulo descreve-se como o modelo apresentado no Capítulo \ref{Met} é realizado como um programa sequencial, no qual um único processador percorre toda a rede a cada passo de tempo.

\section{Estrutura de dados}

Cada indivíduo da população é representado por uma estrutura que reúne o seu estado de saúde e os atributos necessários à evolução da doença, e a rede é representada por uma matriz dessas estruturas. A Listagem \ref{lst:struct} apresenta os principais campos dessa estrutura.

\begin{lstlisting}[caption={Estrutura que representa cada indivíduo (agente).}, label={lst:struct}]
struct Individual {
    int Health;        // estado de saúde atual
    int Swap;          // próximo estado (atualização síncrona)
    int AgeYears;      // idade em anos
    int AgeDays;       // idade em dias
    int AgeDeathYears; // idade de morte natural (anos)
    int AgeDeathDays;  // idade de morte natural (dias)
    int StateTime;     // duração sorteada para o estado atual
    int TimeOnState;   // tempo já decorrido no estado atual
    int Days;          // dias do indivíduo na simulação
    int Isolation;     // indicador de isolamento
    int Exponent;      // tentativas de infecção no passo
    int Checked;       // marca para evitar dupla infecção
} Person[L+2][L+2];
\end{lstlisting}

O campo \texttt{Health} armazena o estado de saúde atual do indivíduo, enquanto o campo \texttt{Swap} armazena o estado para o qual ele transitará, mecanismo detalhado adiante. Os campos de idade guardam a idade atual e a idade de morte natural do indivíduo, e os campos \texttt{StateTime} e \texttt{TimeOnState} guardam, respectivamente, a duração sorteada para o estado atual e o tempo já decorrido nesse estado, de modo que a transição ocorre quando \texttt{TimeOnState} atinge \texttt{StateTime}. Por fim, os campos \texttt{Exponent} e \texttt{Checked} são utilizados no controle do contágio, evitando que um mesmo suscetível seja infectado mais de uma vez no mesmo passo.

A rede é declarada como uma matriz \texttt{Person[L+2][L+2]}, na qual as duas linhas e as duas colunas adicionais formam uma borda que armazena cópias das extremidades opostas da rede, implementando as condições de contorno periódicas descritas no Capítulo \ref{Met}; dessa forma, ao verificar a vizinhança de um indivíduo da borda, o programa encontra automaticamente os vizinhos da borda oposta.

A atualização da rede é feita de forma síncrona por meio de um duplo-buffer, no qual os campos \texttt{Health} e \texttt{Swap} desempenham papéis complementares. Durante a varredura da rede, as funções de estado leem o estado atual de cada indivíduo em \texttt{Health} e escrevem o seu próximo estado em \texttt{Swap}, sem alterar \texttt{Health}; somente ao final do passo o conteúdo de \texttt{Swap} é copiado para \texttt{Health}. Assim, garante-se que a atualização de um indivíduo não interfira no cálculo dos demais no mesmo passo, conforme ilustra a Figura \ref{fig:duplobuffer}.

\begin{figure}[H]
\centering
\caption{Duplo-buffer da atualização síncrona.}
\label{fig:duplobuffer}
\begin{tikzpicture}[font=\small, >=Stealth]
  \node[draw, thick, fill=blue!10, minimum width=4.8cm, minimum height=0.9cm] (health) at (0,1.5) {\texttt{Health} --- estado atual};
  \node[draw, thick, fill=green!12, minimum width=4.8cm, minimum height=0.9cm] (swap) at (0,-1.5) {\texttt{Swap} --- próximo estado};
  \node[draw, thick, rounded corners, fill=gray!10, align=center, text width=3.1cm, minimum height=1.4cm] (func) at (6.4,0) {funções de estado\\ (\texttt{Sfunc}, \texttt{IPfunc}, \dots)};
  \draw[->, thick] (health.east) -- (func.north west) node[midway, above, sloped, font=\scriptsize]{lê};
  \draw[->, thick] (func.south west) -- (swap.east) node[midway, below, sloped, font=\scriptsize]{escreve};
  \draw[->, very thick, dashed] (swap.west) -- ++(-1.8,0) |- (health.west) node[pos=0.5, left, align=center, font=\scriptsize]{atualização\\ síncrona};
\end{tikzpicture}
\end{figure}

Durante a varredura, as funções de estado leem o estado atual em \texttt{Health} e escrevem o próximo estado em \texttt{Swap}; ao final do passo, \texttt{Swap} é copiado para \texttt{Health}, de modo que todas as transições de um passo são aplicadas simultaneamente.

\section{Organização do programa e geração de números aleatórios}

O programa é organizado de forma modular, no qual o arquivo principal contém o laço de simulação e inclui diversos arquivos auxiliares, cada um responsável por uma parte do modelo: a inicialização da população, os parâmetros de cada cidade, as probabilidades de morte natural e de recuperação por faixa etária, os dois mecanismos de contágio e as funções de estado. Cada estado de saúde possui a sua própria função, o que mantém o código legível e preserva a correspondência direta entre o modelo e a sua implementação.

Como o modelo é estocástico, a geração de números aleatórios é parte essencial da implementação. Utiliza-se um gerador congruencial linear multiplicativo, apresentado na Listagem \ref{lst:rng}, no qual o estado do gerador é multiplicado por uma constante a cada chamada e normalizado para o intervalo $[0,1]$. A cada simulação o gerador é inicializado com uma semente distinta, derivada do número da simulação, de modo que cada simulação produz uma sequência própria de números pseudoaleatórios e, ao mesmo tempo, os resultados permanecem reprodutíveis.

\begin{lstlisting}[caption={Gerador congruencial linear multiplicativo.}, label={lst:rng}]
double aleat() {
    R *= mult;                // mult = 888121
    rn = (double) R / MAXNUM; // número em [0, 1]
}
\end{lstlisting}

\section{Inicialização da população}

Antes do início do laço de simulação, a função de inicialização prepara a população de acordo com o descrito no Capítulo \ref{Met}. Toda a rede é colocada no estado suscetível e, a cada indivíduo, é atribuída uma idade sorteada a partir da estrutura etária brasileira, bem como uma idade de morte natural obtida por amostragem por rejeição sobre as probabilidades de morte por faixa etária. Em seguida, os 5 casos iniciais são posicionados aleatoriamente na rede no estado pré-sintomático. A Listagem \ref{lst:begin} apresenta a estrutura dessa função.

\begin{lstlisting}[caption={Inicialização da população (\texttt{beginfunc}).}, label={lst:begin}]
void beginfunc() {
    for (i = 1; i <= L; i++)
      for (j = 1; j <= L; j++) {
        Person[i][j].Health = S;            // toda a rede em suscetível

        // sorteia a faixa etária pela estrutura etária e a idade na faixa
        Person[i][j].AgeYears = rn * (MaximumAge - MinimumAge) + MinimumAge;
        aleat();
        Person[i][j].AgeDays  = rn * 365;

        // idade de morte natural por amostragem por rejeição
        do {
            aleat(); Person[i][j].AgeDeathYears = rn * 100;
            aleat();
        } while (rn >= ProbNaturalDeath[Person[i][j].AgeDeathYears]);
        Person[i][j].AgeDeathDays = Person[i][j].AgeDeathYears * 365;
      }

    // distribui os IPini = 5 casos iniciais em posições aleatórias
    for (k = 0; k < IPini; ) {
        aleat(); i = rn * L + 1;
        aleat(); j = rn * L + 1;
        if (Person[i][j].Health == S) {
            Person[i][j].Health = IP;
            aleat();
            Person[i][j].StateTime = rn * (MaxIP - MinIP) + MinIP;
            k++;
        }
    }
}
\end{lstlisting}

\section{Laço de simulação e funções de estado}

O laço principal segue o procedimento descrito no Algoritmo \ref{alg:simulacao}. Para cada simulação, a população é inicializada e, em seguida, a simulação avança dia a dia: a cada passo, as condições de contorno periódicas são aplicadas e a rede é percorrida sequencialmente, de modo que, para cada indivíduo, é chamada a função correspondente ao seu estado de saúde atual. Concluída a varredura, a rede é atualizada. A Listagem \ref{lst:loop} apresenta esse laço, no qual \texttt{estado\_IS}($\cdot$) abrevia o teste dos estados sintomáticos (IA, $S_l$, $S_m$ e $S_s$).

\begin{lstlisting}[caption={Laço principal: contorno periódico, varredura e atualização.}, label={lst:loop}]
for (sim_time = 0; sim_time <= DAYS; sim_time++) {
    // condições de contorno periódicas (bordas e cantos)
    for (i = 1; i <= L; i++) {
        Person[0][i].Health   = Person[L][i].Health;
        Person[L+1][i].Health = Person[1][i].Health;
        Person[i][0].Health   = Person[i][L].Health;
        Person[i][L+1].Health = Person[i][1].Health;
    }

    // varredura da rede: despacho por estado de saúde
    for (i = 1; i <= L; i++)
      for (j = 1; j <= L; j++)
        if      (Person[i][j].Health == S)   Sfunc(i, j);
        else if (Person[i][j].Health == E)   Efunc(i, j);
        else if (Person[i][j].Health == IP)  IPfunc(i, j);
        else if (estado_IS(Person[i][j].Health)) ISfunc(i, j);
        else if (Person[i][j].Health == H)   Hfunc(i, j);
        else if (Person[i][j].Health == ICU) ICUfunc(i, j);

    Updatefunc();   // atualização síncrona
}
\end{lstlisting}

Cada função de estado segue uma estrutura semelhante: incrementa o tempo decorrido no estado, verifica se o indivíduo sofreu morte natural e, caso contrário, decide a sua permanência ou transição. Enquanto o tempo decorrido não atinge a duração sorteada, o indivíduo permanece no mesmo estado; ao atingi-la, a função sorteia o próximo estado de acordo com as probabilidades de transição e, quando aplicável, com a probabilidade de recuperação correspondente à sua faixa etária, escrevendo o resultado no campo \texttt{Swap}. As funções dos estados infecciosos que ainda circulam na população acrescentam, enquanto o indivíduo permanece infeccioso, a chamada ao segundo mecanismo de contágio, descrito a seguir. A Listagem \ref{lst:ipfunc} ilustra essa estrutura para o estado pré-sintomático.

\begin{lstlisting}[caption={Função de estado do pré-sintomático (\texttt{IPfunc}).}, label={lst:ipfunc}]
void IPfunc(int i, int j) {
    Person[i][j].TimeOnState++;

    if (Person[i][j].Days >= Person[i][j].AgeDeathDays) {   // morte natural
        Person[i][j].Swap = Dead;
        New_Dead++;
    }
    else if (Person[i][j].TimeOnState >= Person[i][j].StateTime) {
        // duração do estado terminou: sorteia o próximo estado
        aleat();
        if (rn < ProbIPtoIA) {
            Person[i][j].Swap = IA;
            aleat();
            Person[i][j].StateTime = rn * (MaxIA - MinIA) + MinIA;
        } else {
            // sorteia ISLight / ISModerate / ISSevere conforme as probabilidades
            // ...
        }
    }
    else {                            // ainda infeccioso
        Neighborsinfectedfunc(i, j);  // segundo mecanismo de contágio
        Person[i][j].Swap = IP;
    }
}
\end{lstlisting}

As funções dos estados sintomático severo, hospitalizado e de UTI consideram ainda a disponibilidade de leitos: um indivíduo em estado severo só é hospitalizado se houver leito de hospital disponível, e a internação em UTI depende da disponibilidade de leitos de UTI. O número de leitos disponíveis em cada cidade desconta a ocupação inicial por outras doenças, conforme descrito no Capítulo \ref{Met}, e é controlado por contadores atualizados ao longo da simulação.

\section{Implementação dos mecanismos de contágio}

O primeiro mecanismo de contágio é realizado pela função associada ao estado suscetível, que percorre a vizinhança local do indivíduo --- de Moore para as cidades de alta densidade e de Von Neumann para as de baixa densidade --- e os seus contatos aleatórios, contando o número $n$ de indivíduos infectados encontrados. A partir desse número, calcula-se a probabilidade de contágio $P = 1 - (1 - \beta)^n$ e, com um número aleatório, decide-se se o suscetível é infectado, caso em que ele transita para o estado exposto. A Listagem \ref{lst:neighbors} apresenta esse mecanismo, no qual \texttt{infeccioso}($\cdot$) denota o teste de o indivíduo estar em um estado infeccioso.

\begin{lstlisting}[caption={Primeiro mecanismo: o suscetível busca infectados (\texttt{Neighborsfunc}).}, label={lst:neighbors}]
int Neighborsfunc(int i, int j) {
    int KI = 0;            // número de infectados encontrados
    Contagion = 0;

    if (Person[i][j].Isolation == IsolationNo) {   // contatos aleatórios
        aleat();
        RandomContacts = rn * (MaxRandomContacts - MinRandomContacts) + MinRandomContacts;
        for (mute = 0; mute < RandomContacts; mute++) {
            // sorteia um contato (Randomi, Randomj)
            if (infeccioso(Randomi, Randomj)) KI++;
        }
    }

#if (Density == HIGH)              // vizinhança de Moore (8 vizinhos)
    for (k = -1; k <= 1; k++)
      for (l = -1; l <= 1; l++)
        if (infeccioso(i + k, j + l)) KI++;
#else                              // vizinhança de Von Neumann (4 vizinhos)
    if (infeccioso(i-1,j) || infeccioso(i+1,j) ||
        infeccioso(i,j-1) || infeccioso(i,j+1)) KI++;
#endif

    if (KI > 0) {                  // P = 1 - (1 - Beta)^KI
        ProbContagion = 1.0 - pow(1.0 - Beta, (double) KI);
        aleat();
        if (rn <= ProbContagion) Contagion = 1;
    }
    return Contagion;
}
\end{lstlisting}

O segundo mecanismo é realizado pelos indivíduos infecciosos que ainda circulam na população, que procuram indivíduos suscetíveis entre seus contatos aleatórios e tentam infectá-los. Para garantir a consistência com a atualização síncrona e evitar que um mesmo suscetível seja infectado duas vezes no mesmo passo, cada suscetível dispõe dos campos \texttt{Exponent} e \texttt{Checked}: o primeiro contabiliza as tentativas de infecção dirigidas àquele indivíduo no passo, enquanto o segundo marca o indivíduo já processado, de modo que a infecção é aplicada uma única vez por passo, independentemente de ter origem na busca do próprio suscetível ou na de um infectado. A Listagem \ref{lst:neighborsinf} apresenta esse mecanismo.

\begin{lstlisting}[caption={Segundo mecanismo: o infectado busca suscetíveis (\texttt{Neighborsinfectedfunc}).}, label={lst:neighborsinf}]
void Neighborsinfectedfunc(int i, int j) {
    if (Person[i][j].Isolation == IsolationNo) {
        aleat();
        RandomContacts = rn * (MaxRandomContacts - MinRandomContacts) + MinRandomContacts;
        for (mute = 0; mute < RandomContacts; mute++) {
            // sorteia um suscetível-alvo (Randomi, Randomj)
            if (Person[Randomi][Randomj].Health == S)
                Person[Randomi][Randomj].Exponent++;

            if (Person[Randomi][Randomj].Checked == 0 &&
                Person[Randomi][Randomj].Exponent == 1) {
                ProbContagion = 1.0 - pow(1.0 - Beta, (double) Person[Randomi][Randomj].Exponent);
                aleat();
                if (rn <= ProbContagion) {
                    Person[Randomi][Randomj].Checked = 1;
                    Person[Randomi][Randomj].Swap = E;
                    Person[Randomi][Randomj].TimeOnState = 0;
                    aleat();
                    Person[Randomi][Randomj].StateTime = rn * (MaxLatency - MinLatency) + MinLatency;
                }
            }
        }
    }
}
\end{lstlisting}

\section{Atualização síncrona e geração das saídas}

Ao final de cada passo de tempo, a função de atualização copia o campo \texttt{Swap} para o campo \texttt{Health} de todos os indivíduos, efetivando de forma síncrona as transições calculadas durante a varredura, e incrementa os contadores de idade e de tempo de simulação. Em seguida, contabiliza-se quantos indivíduos se encontram em cada estado de saúde, e os indivíduos que atingiram a sua idade de morte natural são removidos e substituídos por novos suscetíveis, com idade sorteada novamente a partir da estrutura etária, de modo que o tamanho total da população permanece constante. A Listagem \ref{lst:update} apresenta o trecho central dessa atualização.

\begin{lstlisting}[caption={Atualização síncrona e substituição dos mortos (\texttt{Updatefunc}).}, label={lst:update}]
for (i = 1; i <= L; i++)
  for (j = 1; j <= L; j++) {
    Person[i][j].Health = Person[i][j].Swap;   // atualização síncrona
    Person[i][j].Exponent = 0;
    Person[i][j].Checked  = 0;
    Person[i][j].AgeDays++;
    Person[i][j].Days++;

    if (Person[i][j].Health == Dead || Person[i][j].Health == DeadCovid) {
        Person[i][j].Health = S;                // novo suscetível
        aleat();
        Person[i][j].AgeYears = rn * 100;       // idade reamostrada
        // ... sorteia a idade de morte natural (amostragem por rejeição)
    }
  }
\end{lstlisting}

Os resultados são registrados em dois níveis. Para cada simulação, a prevalência e a incidência de cada estado de saúde são gravadas, a cada dia, em arquivos próprios. Ao final de todas as simulações, calcula-se a média das curvas sobre as \texttt{MAXSIM} simulações, obtendo-se, pelo método de Monte Carlo, as curvas médias de prevalência e de incidência que representam a evolução esperada da epidemia.
